
\section{The WildIFEval Replication}
\label{sec:appendix:wild}

\subsection{The design grid}
\label{sec:appendix:wildgrid}

\begin{table*}[t]
  \centering\small
  \begin{tabular}{@{}rrrlrrrrrrr@{}}
    \toprule
    & & & & & \multicolumn{3}{c}{\textsc{ssr}} & & & \\
    \cmidrule(lr){6-8}
    $N$ & $K$ & $D$ & Arm & $n$ & provided & withheld & all & $|D_t|$ & $|D_T|$ & trunc. \\
    \midrule
    5 & 1 & 0 & \textsc{control} & 32 & 0.547 & 0.531 & 0.544 & 5.500 & 6.594 & 0.219 \\
    5 & 1 & 0 & \textsc{placebo} & 32 & 0.578 & 0.500 & 0.562 & 5.250 & 5.781 & 0.188 \\
    5 & 1 & 0 & \textsc{treatment} & 32 & 0.633 & 0.531 & 0.613 & 5.750 & 6.250 & 0.281 \\
    \addlinespace
    5 & 1 & 16 & \textsc{control} & 32 & 0.531 & 0.438 & 0.512 & 14.750 & 14.406 & 0.312 \\
    5 & 1 & 16 & \textsc{placebo} & 32 & 0.523 & 0.500 & 0.519 & 14.406 & 12.969 & 0.312 \\
    5 & 1 & 16 & \textsc{treatment} & 32 & 0.656 & 0.500 & 0.625 & 6.500 & 7.031 & 0.250 \\
    \addlinespace
    6 & 1 & 0 & \textsc{control} & 32 & 0.688 & 0.531 & 0.661 & 6.656 & 6.969 & 0.375 \\
    6 & 1 & 0 & \textsc{placebo} & 32 & 0.650 & 0.500 & 0.625 & 6.312 & 6.906 & 0.344 \\
    6 & 1 & 0 & \textsc{treatment} & 32 & 0.575 & 0.500 & 0.562 & 6.562 & 6.594 & 0.438 \\
    \addlinespace
    6 & 1 & 16 & \textsc{control} & 32 & 0.525 & 0.344 & 0.495 & 14.812 & 14.844 & 0.344 \\
    6 & 1 & 16 & \textsc{placebo} & 32 & 0.569 & 0.406 & 0.542 & 14.969 & 13.688 & 0.344 \\
    6 & 1 & 16 & \textsc{treatment} & 32 & 0.631 & 0.531 & 0.615 & 7.125 & 7.562 & 0.406 \\
    \bottomrule
  \end{tabular}
  \caption{
    The arm effect lives where the distractors are. Every design cell of \S\ref{sec:wild} at the
    final maintenance round: $N$ hidden constraints, $K$ of them withheld from the seeded prompt,
    $D$ distractors seeded alongside the $N{-}K$ true instructions, one row per arm. The
    \textsc{ssr} columns give the per-constraint satisfaction rate over the provided $N{-}K$
    (interference), the withheld $K$ (recovery), and all $N$ (the headline); $|D_t|$ is the prompt
    the executor saw, $|D_T|$ the prompt after the last round, and trunc.\ the share of responses
    that hit the executor's token cap, tabulated beside the rates it could otherwise explain. At
    $D{=}16$ \textsc{treatment} holds all-$N$ \textsc{ssr} at its $D{=}0$ level while
    \textsc{control} falls in both strata. Size columns are descriptive only: this run's seeded
    comment states its own falsification, so pruning a distractor needs no inference from the
    maintainer, and no size claim in this paper cites this run.
  }
  \label{tab:wildgrid}
\end{table*}


\autoref{sec:wild} reports four contrasts, and \autoref{tab:wildgrid} reports every cell they are
computed from, so a reader can check the headline against the design rather than against one
favorable comparison. The grid is the same computation the analysis notebook renders, restricted
to this suite.

\subsection{Scoring}
\label{sec:appendix:wildscoring}

WildIFEval ships human-written constraint decompositions and no code verifiers,
so satisfaction here is an LLM judge rather than \autoref{tab:verifiers}'s programmatic ones. The judge
is \texttt{claude-haiku-4-5}, prompted with exactly one constraint and one response and asked for a
boolean --- no prompt $D_t$, no arm label, no history, no sibling constraints --- so nothing in its input
distinguishes the arms, and two arms producing the same response are scored identically. Both arms are
judged by the same model in the same call shape, so judge error is common to the contrast rather than
differential across it. Responses are stored verbatim, so the suite can be re-judged without
re-generating them, which is what makes the next subsection cost no generation.

\subsection{Judge robustness}
\label{sec:appendix:judges}

\begin{table*}[t]
  \centering\footnotesize
  \begin{tabular}{@{}lrrr@{}}
    \toprule
    & \multicolumn{2}{c}{Contrast under each judge} & \textsc{did} \\
    \cmidrule(lr){2-3}
    Criterion & \texttt{claude-haiku-4-5} & \texttt{gpt-5.6-luna} & \shortstack[r]{\texttt{claude-haiku-4-5}\\$-$\,\texttt{gpt-5.6-luna}} \\
    \midrule
    (i) interference & $\mathbf{-0.241}$ [$-0.334$, $-0.149$] & $\mathbf{-0.197}$ [$-0.287$, $-0.108$] & --- \\
    (ii) comments buy IF & $\mathbf{+0.116}$ [$+0.051$, $+0.183$] & $\mathbf{+0.078}$ [$+0.006$, $+0.149$] & $+0.038$ [$-0.019$, $+0.096$] \\
    (v) not just recovery & $\mathbf{+0.116}$ [$+0.048$, $+0.185$] & $\mathbf{+0.083}$ [$+0.005$, $+0.159$] & $+0.033$ [$-0.031$, $+0.095$] \\
    (aux) recovery & $\mathbf{+0.125}$ [$+0.016$, $+0.250$] & $+0.062$ [$-0.062$, $+0.188$] & $+0.062$ [$-0.062$, $+0.203$] \\
    (iii) comments buy size & $\mathbf{-7.328}$ [$-8.641$, $-5.953$] & $\mathbf{-7.328}$ [$-8.641$, $-5.953$] & $+0.000$ [$+0.000$, $+0.000$] \\
    (iv) placebo null (SSR) & $+0.027$ [$-0.044$, $+0.101$] & $+0.033$ [$-0.036$, $+0.101$] & $-0.006$ [$-0.070$, $+0.061$] \\
    (iv) placebo null (size) & $-1.297$ [$-2.656$, $+0.031$] & $-1.297$ [$-2.656$, $+0.031$] & $+0.000$ [$+0.000$, $+0.000$] \\
    \bottomrule
  \end{tabular}
  \caption{
    Every contrast \S\ref{sec:wild} reports survives an independent scorer, and the two judges'
    effect sizes are not distinguishable. Rows are the pre-registered criteria of
    \apxref{sec:appendix:wildgrid}; columns 2--3 give the paired
    \textsc{treatment}$-$\textsc{control} difference under each judge (\textsc{placebo}$-$%
    \textsc{control} for the placebo rows, and $D{=}16$ vs.\ $D{=}0$ within \textsc{control}
    for interference), with its 95\% percentile bootstrap interval over $W{=}64$ worlds;
    \textbf{bold} marks a point estimate whose interval excludes zero. \texttt{gpt-5.6-luna} matches
    \texttt{claude-haiku-4-5}'s capability tier, so a disagreement between them is a vendor difference rather
    than a capability one, and it re-scores the \emph{same} stored responses through the same
    prompt, schema, token cap and $\text{temperature}{=}0$: only the model differs. Column 4 is the
    per-world difference-in-differences $(\text{T}-\text{C})|_a - (\text{T}-\text{C})|_b$, for $a$
    and $b$ the judges of columns 2 and 3, over those same worlds --- the estimate that says whether
    two judges' effect sizes differ, which two overlapping per-judge intervals do not. It covers zero on every row, so the spread between the
    two point estimates has no evidential support as a judge difference. At $W{=}64$ that
    is absence of evidence and not equivalence: the interval on (ii) reaches $0.096$. Rates
    are per-constraint satisfaction; the two size rows are counts of directives, and are the
    built-in control --- no judge sees $|D_T|$, so their columns must agree exactly and their
    \textsc{did} must be $0$.
  }
  \label{tab:judges}
\end{table*}

\begin{table*}[t]
  \centering\footnotesize
  \begin{tabular}{@{}lrrrrrr@{}}
    \toprule
    & & & & \multicolumn{2}{c}{Pass rate} & \\
    \cmidrule(lr){5-6}
    Slice & $n$ & Agree & $\kappa$ [95\% CI] & \texttt{claude-haiku-4-5} & \texttt{gpt-5.6-luna} & Lenience \\
    \midrule
    all arms & 6,336 & 0.821 & 0.639 [0.620, 0.657] & 0.545 & 0.547 & $+0.002$ \\
    \addlinespace
    \textsc{control} & 2,112 & 0.816 & 0.631 [0.598, 0.664] & 0.539 & 0.534 & $-0.006$ \\
    \textsc{treatment} & 2,112 & 0.819 & 0.632 [0.598, 0.665] & 0.565 & 0.563 & $-0.002$ \\
    \textsc{placebo} & 2,112 & 0.827 & 0.653 [0.618, 0.684] & 0.529 & 0.543 & $+0.014$ \\
    \bottomrule
  \end{tabular}
  \caption{
    Judge agreement is substantial and does not vary with the arm, so judge noise cannot have
    manufactured the arm gap in \autoref{tab:judges}. Judge $a = $ \texttt{claude-haiku-4-5} against judge
    $b = $ \texttt{gpt-5.6-luna} on all 6,336 verdicts, pooled and per arm: raw agreement, Cohen's
    $\kappa$ with its 95\% bootstrap interval, each judge's own marginal pass rate, and lenience,
    the signed disagreement $(b\text{ only} - a\text{ only})/n$ --- positive means $b$ passes what
    $a$ fails. $\kappa$ is shown beside the marginals throughout because it collapses toward zero
    whenever one label is rare however well two raters agree; here the marginals are near-equal, so
    $\kappa$ is not a prevalence artifact. Read the per-arm rows first: uniform lenience compresses
    every arm toward the ceiling and shrinks all gaps mechanically, whereas lenience that varied
    \emph{by arm} would mean the judge interacts with the treatment. It does not vary --- $\kappa$
    spans 0.631--0.653 across the three arms. The judges disagree on
    17.9\% of individual verdicts while their marginal pass rates differ by
    0.002, so the disagreements cancel rather than accumulate: same threshold,
    different items. Judge~1 decoded at the vendor default and the re-judge at
    $\text{temperature}{=}0$, so these figures absorb judge~1's sampling variance and bound the
    agreement from below.
  }
  \label{tab:judgeagreement}
\end{table*}


No contrast in \autoref{sec:wild} is a property of the model that graded it. We re-scored all
\JudgeVerdicts{} stored verdicts under \JudgeAlt{}, a second judge matched to judge~1's capability
tier so that a disagreement between them is a vendor difference rather than a capability one,
holding the prompt, schema, token cap and $\text{temperature}{=}0$ fixed so that only the model
varies. \autoref{tab:judges} is the result: every criterion judge~1 passes keeps its pre-registered
sign under the second judge, interference and the size criteria keep intervals excluding zero under
both, and the placebo stays null under both. The two judges also agree on the effect's \emph{size}
as far as this design can tell. Their per-world difference-in-differences on the
instruction-following criterion is \JudgeDidPp{} (95\% CI: \JudgeDidCI{}), so the spread
between their point estimates (\WildDeltaSatPp{} and \JudgeAltDeltaSatPp{}) has no
evidential support as a judge difference. At $W = \WildWorlds{}$ that is absence of evidence rather
than equivalence: the interval admits a judge difference as large as \JudgeDidBound{}, and
closing that needs more worlds rather than more judges.

\subsection{Judge agreement, and the scope of the effect}
\label{sec:appendix:judgeagreement}

\autoref{tab:judgeagreement} reports agreement, which we record but do not gate on: two judges can
disagree about absolute difficulty and still rank the arms identically, and it is the ranking
\autoref{sec:wild} quotes. Judge~1 and \JudgeAlt{} agree on \JudgeAgreePct{} of individual
verdicts at $\kappa = \JudgeKappa{}$ ($\JudgeKappaCI$), and --- the load-bearing part --- $\kappa$
does not vary with the arm, so judge noise cannot have produced the arm gap. What we do not measure
is judge \emph{accuracy}: neither model is ground truth, so a disagreement never says which one is
right, and every \autoref{sec:wild} rate remains a rate under a named judge. Licensing an accuracy
claim would take a human-annotated subsample, which we did not run.

The effect is a recovery from seeded noise rather than a general gain. At $D = 0$ the arms separate
in neither direction consistently: the commented arm sits above its control at $N = 5$ and below
it at $N = 6$. The commented arm holds satisfaction while the control falls in both strata only at
$D = 16$, where the prompt carries distractors to prune. A prompt carrying no extraneous
instructions is outside the scope of this experiment.

